\documentclass[12pt, a4paper]{article}

% ---------------------------------------------------------------
% Paquetes
% ---------------------------------------------------------------
\usepackage[spanish]{babel}
\usepackage[utf8]{inputenc}
\usepackage[T1]{fontenc}
\usepackage[margin=2.5cm]{geometry}
\usepackage{hyperref}
\usepackage{xcolor}
\usepackage{listings}
\usepackage{enumitem}
\usepackage{fancyhdr}
\usepackage{titlesec}
\usepackage[skins, breakable]{tcolorbox}
\usepackage{booktabs}
\usepackage{graphicx}
\usepackage{parskip}
\usepackage{array}
\usepackage{tabularx}

% ---------------------------------------------------------------
% Colores
% ---------------------------------------------------------------
\definecolor{uccazul}{RGB}{0, 70, 127}
\definecolor{uccgris}{RGB}{80, 80, 80}
\definecolor{uccverde}{RGB}{22, 163, 74}
\definecolor{alertanaranja}{RGB}{230, 126, 34}
\definecolor{alertarojo}{RGB}{220, 38, 38}
\definecolor{codegris}{RGB}{245, 245, 245}
\definecolor{amarillo}{RGB}{202, 138, 4}

% ---------------------------------------------------------------
% Cajas
% ---------------------------------------------------------------
\newtcolorbox{aviso}{
  colback=alertanaranja!10, colframe=alertanaranja,
  title={\bfseries Importante}, fonttitle=\bfseries,
  breakable
}
\newtcolorbox{nota}{
  colback=uccazul!8, colframe=uccazul,
  title={\bfseries Nota}, fonttitle=\bfseries,
  breakable
}
\newtcolorbox{tip}{
  colback=uccverde!8, colframe=uccverde,
  title={\bfseries Consejo}, fonttitle=\bfseries,
  breakable
}
\newtcolorbox{peligro}{
  colback=alertarojo!8, colframe=alertarojo,
  title={\bfseries Atención}, fonttitle=\bfseries,
  breakable
}

% ---------------------------------------------------------------
% Código
% ---------------------------------------------------------------
\lstset{
  backgroundcolor=\color{codegris},
  basicstyle=\ttfamily\small,
  breaklines=true,
  frame=single,
  framesep=5pt,
  rulecolor=\color{gray!40},
  xleftmargin=10pt,
  xrightmargin=10pt,
}

% ---------------------------------------------------------------
% Encabezado y pie
% ---------------------------------------------------------------
\pagestyle{fancy}
\fancyhf{}
\fancyhead[L]{\color{uccgris}\small Sistema de Control de Acceso UCC}
\fancyhead[R]{\color{uccgris}\small Manual de Uso}
\fancyfoot[C]{\color{uccgris}\small \thepage}
\renewcommand{\headrulewidth}{0.4pt}

% ---------------------------------------------------------------
% Títulos
% ---------------------------------------------------------------
\titleformat{\section}
  {\Large\bfseries\color{uccazul}}{\thesection.}{0.5em}{}[\titlerule]
\titleformat{\subsection}
  {\large\bfseries\color{uccazul!80}}{\thesubsection.}{0.5em}{}
\titleformat{\subsubsection}
  {\normalsize\bfseries\color{uccgris}}{\thesubsubsection.}{0.5em}{}

% ---------------------------------------------------------------
% Hipervínculos
% ---------------------------------------------------------------
\hypersetup{
  colorlinks=true,
  linkcolor=uccazul,
  urlcolor=uccazul,
  pdftitle={Manual de Uso - UCC Control de Acceso},
  pdfauthor={Corporacion Universitaria Comfacauca}
}

% ===============================================================
\begin{document}
% ===============================================================

% === Portada ===
\begin{titlepage}
  \centering
  \vspace*{2.5cm}
  {\color{uccazul}\rule{\textwidth}{2.5pt}}\\[1em]
  {\Huge\bfseries\color{uccazul} Sistema de Control de Acceso}\\[0.6em]
  {\LARGE\color{uccgris} Corporación Universitaria Comfacauca}\\[0.6em]
  {\color{uccazul}\rule{\textwidth}{2.5pt}}\\[2em]

  {\Large\bfseries Manual de Uso del Sistema}\\[0.5em]
  {\large\color{uccgris} Versión 1.0 — 2026}

  \vfill

  \begin{tcolorbox}[colback=uccazul!5, colframe=uccazul!30, width=0.85\textwidth]
    \centering
    \small\color{uccgris}
    Este manual describe cómo usar el sistema tanto desde el
    \textbf{portal de usuarios} (estudiantes, empleados y contratistas)
    como desde el \textbf{panel de administración}.\\[0.4em]
    No se requieren conocimientos técnicos previos.
  \end{tcolorbox}

  \vspace{2cm}
  {\color{uccgris}\small Corporación Universitaria Comfacauca, 2026}
\end{titlepage}

% === Tabla de contenidos ===
\tableofcontents
\newpage

% ===============================================================
\section{¿Qué es este sistema?}
% ===============================================================

El Sistema de Control de Acceso UCC es una aplicación web que permite
registrar y verificar quién entra a las instalaciones de la universidad.
Funciona en la red WiFi del campus: cualquier persona con un celular o
computador conectado a esa red puede usarlo.

El sistema tiene dos interfaces distintas con propósitos diferentes:

\begin{center}
\renewcommand{\arraystretch}{1.5}
\begin{tabular}{@{} p{3.5cm} p{5.5cm} p{4.5cm} @{}}
  \toprule
  \textbf{Interfaz} & \textbf{Para quién} & \textbf{Cómo acceder} \\
  \midrule
  Portal de usuario &
    Estudiantes, empleados y contratistas que registran su ingreso &
    Código QR o \texttt{http://[IP]/login} \\[0.4em]
  Panel de administración &
    Personal administrativo que gestiona el sistema &
    \texttt{http://[IP]/admin-login} \\
  \bottomrule
\end{tabular}
\end{center}

\begin{nota}
  En este manual, \texttt{[IP]} representa la dirección de red del equipo
  donde está instalado el sistema (por ejemplo, \texttt{192.168.44.216}).
  El administrador del sistema debe informarte cuál es esa dirección.
\end{nota}

% ===============================================================
\section{Portal de usuario}
\label{sec:portal}
% ===============================================================

El portal de usuario es la pantalla que verán los estudiantes,
empleados y contratistas cuando escaneen el código QR en la entrada
de la universidad.

\subsection{¿Cómo ingresar?}

Hay dos formas de llegar al portal:

\begin{enumerate}
  \item \textbf{Código QR}: Escanea el código QR publicado en la entrada
        con la cámara del celular. Esto abre automáticamente el portal
        en el navegador.

  \item \textbf{URL directa}: Desde cualquier navegador (Chrome, Firefox,
        Edge, Safari) conectado a la red WiFi del campus, escribe la
        dirección del portal:
        \begin{lstlisting}
http://[IP]/login
        \end{lstlisting}
\end{enumerate}

\subsection{Registrar el ingreso}

Al abrir el portal, verás una pantalla con el logo de la UCC y un campo
para ingresar tu número de identificación.

\begin{enumerate}
  \item Escribe tu \textbf{ID institucional} (el número que te asignó la
        universidad al momento de tu vinculación).
  \item Presiona el botón \textbf{Ingresar} o la tecla Enter.
  \item El sistema verifica tu información y muestra el resultado.
\end{enumerate}

\subsubsection{Resultado exitoso}

Si tu ID es válido y tu cuenta está activa, verás una pantalla de
confirmación con:
\begin{itemize}
  \item Tu nombre completo.
  \item Tu rol (estudiante, empleado o contratista).
  \item La hora exacta del registro.
\end{itemize}

\subsubsection{Resultado con error}

\begin{center}
\renewcommand{\arraystretch}{1.4}
\begin{tabular}{@{} p{4.5cm} p{8.5cm} @{}}
  \toprule
  \textbf{Mensaje} & \textbf{Qué significa} \\
  \midrule
  Usuario no encontrado &
    El ID no existe en el sistema. Verifica que lo hayas escrito
    correctamente o contacta a la administración. \\[0.3em]
  Acceso bloqueado &
    Tu cuenta ha sido suspendida temporalmente. Debes acercarte
    a la oficina de atención para resolver la situación. \\[0.3em]
  Error de conexión &
    El sistema no está disponible en este momento. Avisa al
    personal administrativo. \\
  \bottomrule
\end{tabular}
\end{center}

\subsection{Reglamento TIP}

En la parte inferior del portal hay un enlace para descargar el
\textbf{Reglamento TIP} (Resolución Rectoral N° 258 de 2012). Este
documento describe las normas de convivencia y los motivos por los
cuales una cuenta puede ser bloqueada.

\begin{aviso}
  El portal de usuario está diseñado para ser usado desde el celular,
  pero también funciona desde cualquier computador. No requiere
  instalar ninguna aplicación.
\end{aviso}

% ===============================================================
\section{Panel de administración}
\label{sec:admin}
% ===============================================================

El panel de administración es la interfaz con la que el personal
encargado gestiona el sistema: revisa estadísticas, administra usuarios,
genera reportes y carga listas de personas al inicio de cada semestre.

\subsection{Acceder al panel}

\begin{enumerate}
  \item Abre un navegador web en cualquier equipo conectado a la red
        del campus.
  \item Escribe la siguiente dirección:
        \begin{lstlisting}
http://[IP]/admin-login
        \end{lstlisting}
        También puedes acceder desde el mismo equipo donde está el
        sistema usando:
        \begin{lstlisting}
http://localhost/admin-login
        \end{lstlisting}
  \item Ingresa tu \textbf{ID institucional} y tu \textbf{contraseña}.
  \item Presiona \textbf{Ingresar}.
\end{enumerate}

\begin{nota}
  Las credenciales del panel de administración son completamente
  independientes de las del portal de usuario. Tener una cuenta de
  usuario no te da acceso al panel de administración.
\end{nota}

\begin{aviso}
  Si es la primera vez que usas el sistema, las credenciales por
  defecto son:\\
  \textbf{ID:} \texttt{000000} \quad \textbf{Contraseña:} \texttt{12345678}\\
  Cámbialas inmediatamente después del primer ingreso.
\end{aviso}

\subsection{Niveles de administrador}

El sistema maneja dos niveles de acceso para administradores:

\begin{center}
\renewcommand{\arraystretch}{1.4}
\begin{tabular}{@{} p{3cm} p{9.5cm} @{}}
  \toprule
  \textbf{Nivel} & \textbf{Permisos} \\
  \midrule
  \textbf{Admin} &
    Puede ver el resumen, gestionar usuarios, cargar CSV,
    ver informes y generar reportes. \\[0.3em]
  \textbf{Superadmin} &
    Todo lo anterior, más la capacidad de crear y eliminar
    otros administradores del sistema. \\
  \bottomrule
\end{tabular}
\end{center}

\subsection{Navegación del panel}

Una vez dentro del panel, verás una barra lateral a la izquierda con
el menú de navegación. Contiene las siguientes secciones:

\begin{center}
\renewcommand{\arraystretch}{1.4}
\begin{tabular}{@{} p{3.5cm} p{9cm} @{}}
  \toprule
  \textbf{Sección} & \textbf{Descripción} \\
  \midrule
  Resumen         & Dashboard principal con estadísticas y gráficas en tiempo real. \\
  Usuarios        & Listado completo de personas registradas con búsqueda y filtros. \\
  Semestre        & Carga masiva de usuarios al inicio de cada período académico. \\
  Informes        & Generación y exportación de reportes de acceso. \\
  Administradores & Gestión de cuentas de administrador (solo superadmin). \\
  \bottomrule
\end{tabular}
\end{center}

En la parte inferior de la barra lateral hay un indicador del
\textbf{estado del sistema} que muestra si la base de datos está
conectada correctamente, y el botón de \textbf{cerrar sesión}.

% ===============================================================
\section{Sección: Resumen}
\label{sec:resumen}
% ===============================================================

La sección de Resumen es la pantalla principal del panel. Muestra
una fotografía del estado actual del sistema con datos actualizados
en tiempo real.

\subsection{Tarjetas de estadísticas}

En la parte superior hay ocho tarjetas con los indicadores principales:

\begin{itemize}
  \item \textbf{Estudiantes}: Total de estudiantes registrados en el semestre activo.
  \item \textbf{Activos}: Personas con acceso habilitado.
  \item \textbf{Bloqueados}: Personas con acceso suspendido.
  \item \textbf{Reportes TIP hoy}: Cantidad de reportes disciplinarios registrados hoy.
  \item \textbf{Empleados}: Total de empleados registrados.
  \item \textbf{Contratistas}: Total de contratistas registrados.
  \item \textbf{Ingresos hoy}: Total de registros de acceso del día actual.
  \item \textbf{Semestre activo}: Nombre del período académico en curso.
\end{itemize}

\subsection{Gráficas}

Debajo de las tarjetas aparecen tres visualizaciones:

\begin{enumerate}
  \item \textbf{Gráfica de dona (Activos vs. Bloqueados)}: Muestra la
        proporción de usuarios activos frente a los bloqueados, con la
        tasa de bloqueo expresada en porcentaje.

  \item \textbf{Reportes TIP últimos 7 días}: Barra vertical que muestra
        cuántos reportes disciplinarios se registraron cada día en la
        última semana.

  \item \textbf{Distribución por programa}: Barra horizontal que muestra
        cuántos estudiantes hay en cada programa académico.
\end{enumerate}

\subsection{Top 5 usuarios con más reportes}

Al final de la sección aparece una lista con los cinco usuarios que
acumulan más reportes TIP en el semestre. Esto permite identificar
rápidamente casos que requieren atención.

\subsection{Botón de actualizar}

Las estadísticas no se actualizan solas de forma continua. Para ver
los datos más recientes, presiona el botón de \textbf{actualizar}
(ícono de flecha circular) ubicado en la esquina superior derecha
de la sección.

% ===============================================================
\section{Sección: Usuarios}
\label{sec:usuarios}
% ===============================================================

La sección de Usuarios muestra el listado completo de todas las personas
registradas en el sistema, con herramientas para buscar, filtrar y
gestionar el acceso de cada una.

\subsection{Buscar un usuario}

En la parte superior hay una barra de búsqueda. Puedes escribir
cualquier combinación de:
\begin{itemize}
  \item Nombre completo o parte de él.
  \item ID institucional.
  \item Rol (estudiante, empleado, contratista).
  \item Información adicional (programa académico, cargo, empresa).
\end{itemize}

El listado se filtra en tiempo real mientras escribes.
Para limpiar la búsqueda, presiona la \textbf{X} al final del campo.

\subsection{Filtrar por estado}

Junto a la barra de búsqueda hay tres botones de filtro rápido:

\begin{itemize}
  \item \textbf{Todos}: Muestra estudiantes, empleados y contratistas.
  \item \textbf{Activo}: Solo personas con acceso habilitado.
  \item \textbf{Bloqueado}: Solo personas con acceso suspendido.
\end{itemize}

\subsection{Información mostrada por usuario}

Cada fila del listado muestra:
\begin{itemize}
  \item Nombre completo e ID institucional.
  \item Rol (con colores diferenciadores: azul para estudiantes,
        morado para empleados, naranja para contratistas).
  \item Estado actual (verde = activo, rojo = bloqueado).
  \item Información adicional (programa, cargo o empresa según corresponda).
  \item Cantidad de reportes TIP acumulados.
\end{itemize}

\subsection{Bloquear o desbloquear un usuario}

Para cambiar el estado de acceso de un usuario:

\begin{enumerate}
  \item Ubica al usuario en el listado (usa la búsqueda si es necesario).
  \item Haz clic en el botón de estado que aparece en la columna derecha
        de esa fila.
        \begin{itemize}
          \item Si el usuario está \textbf{activo}, el botón dice
                \emph{"Bloquear"} en rojo.
          \item Si está \textbf{bloqueado}, el botón dice
                \emph{"Habilitar"} en verde.
        \end{itemize}
  \item Aparecerá una ventana de confirmación con el nombre del usuario
        y la acción que vas a realizar.
  \item Presiona \textbf{Confirmar} para aplicar el cambio o
        \textbf{Cancelar} para volver.
\end{enumerate}

\begin{aviso}
  Esta acción tiene efecto inmediato. Una vez bloqueado, el usuario
  no podrá registrar su ingreso en el portal hasta que sea habilitado
  nuevamente.
\end{aviso}

\begin{tip}
  Si el listado tarda en cargar, usa el botón de \textbf{actualizar}
  que aparece en la esquina superior derecha de esta sección.
\end{tip}

% ===============================================================
\section{Sección: Semestre (Carga masiva)}
\label{sec:semestre}
% ===============================================================

Al inicio de cada semestre académico, se carga la lista actualizada
de estudiantes, empleados y contratistas. Esta sección guía el proceso
paso a paso con un asistente en forma de \emph{stepper} (serie de pasos
numerados en la parte superior de la pantalla).

\subsection{¿Cuándo usar esta sección?}

\begin{itemize}
  \item Al inicio de un nuevo semestre académico.
  \item Cuando llegan nuevas listas de personal (empleados o contratistas).
  \item Cuando se necesita reemplazar la base de usuarios vigente.
\end{itemize}

\begin{peligro}
  La carga de un nuevo semestre \textbf{reemplaza} la lista de usuarios
  anterior. Los registros de acceso (ingresos históricos) no se borran,
  pero los usuarios que no aparezcan en los nuevos archivos CSV
  quedarán sin información de semestre activo.
\end{peligro}

\subsection{Pasos del proceso}

El asistente tiene 6 pasos. Debes completarlos en orden.

\subsubsection{Paso 1 — Semestre}

Define el período académico que vas a cargar.

\begin{enumerate}
  \item Si el semestre ya existe en el sistema (aparece en la lista
        desplegable), selecciónalo.
  \item Si es un semestre nuevo, escribe el nombre (por ejemplo:
        \texttt{2026-2}), la fecha de inicio y la fecha de fin.
  \item Presiona \textbf{Siguiente}.
\end{enumerate}

\subsubsection{Paso 2 — Usuarios}

Sube el archivo CSV con la lista general de usuarios.

\begin{enumerate}
  \item Prepara el archivo CSV con el siguiente formato de columnas:
\end{enumerate}

\begin{lstlisting}
id_institucional, nombre_completo, acceso
100002, Ana Torres, activo
100003, Luis Pérez, activo
\end{lstlisting}

\begin{itemize}
  \item \texttt{id\_institucional}: Número de identificación único.
  \item \texttt{nombre\_completo}: Nombre y apellidos.
  \item \texttt{acceso}: Puede ser \texttt{activo} o \texttt{bloqueado}.
\end{itemize}

\begin{enumerate}
  \setcounter{enumi}{1}
  \item Haz clic en \textbf{Seleccionar archivo} y elige tu CSV.
  \item El sistema muestra una vista previa con los datos detectados.
        Verifica que los nombres de las columnas sean correctos.
  \item Presiona \textbf{Siguiente}.
\end{enumerate}

\subsubsection{Paso 3 — Estudiantes}

Sube el CSV con la información académica de los estudiantes.

\begin{lstlisting}
id_institucional, programa, semestre_num
100002, Ingeniería de Sistemas, 5
100003, Contaduría Pública, 3
\end{lstlisting}

\begin{nota}
  Solo deben aparecer en este archivo los IDs que también estaban en
  el CSV del Paso 2 con rol estudiante. Los IDs desconocidos
  serán ignorados.
\end{nota}

\subsubsection{Paso 4 — Empleados}

Sube el CSV con la información laboral de los empleados.

\begin{lstlisting}
id_institucional, cargo, departamento
200001, Docente, Facultad de Ingeniería
200002, Coordinador, Bienestar Universitario
\end{lstlisting}

\subsubsection{Paso 5 — Contratistas}

Sube el CSV con la información de los contratistas externos.

\begin{lstlisting}
id_institucional, empresa, cargo
300001, Servicios Generales S.A.S, Operario
300002, TechSoluciones Ltda, Técnico
\end{lstlisting}

\subsubsection{Paso 6 — Confirmar}

El sistema muestra un resumen de todos los datos que se van a cargar:
\begin{itemize}
  \item Nombre del semestre seleccionado.
  \item Cantidad de usuarios por archivo.
  \item Alertas sobre IDs duplicados o datos faltantes.
\end{itemize}

Si todo está correcto, presiona \textbf{Confirmar y cargar}.
El sistema procesa los archivos y muestra un mensaje de éxito al terminar.

\begin{tip}
  Puedes usar los archivos de ejemplo incluidos en la carpeta
  \texttt{database\textbackslash{}} del sistema como plantilla para
  preparar tus propios CSV:
  \begin{itemize}
    \item \texttt{test\_estudiantes.csv}
    \item \texttt{test\_empleados.csv}
    \item \texttt{test\_contratistas.csv}
    \item \texttt{ejemplo\_carga\_semestral.csv}
  \end{itemize}
\end{tip}

% ===============================================================
\section{Sección: Informes}
\label{sec:informes}
% ===============================================================

La sección de Informes permite consultar y exportar reportes de
acceso filtrando por diferentes períodos de tiempo.

\subsection{Seleccionar el período}

En la parte superior de la sección hay cuatro opciones de período:

\begin{center}
\renewcommand{\arraystretch}{1.4}
\begin{tabular}{@{} p{2.5cm} p{9.5cm} @{}}
  \toprule
  \textbf{Período} & \textbf{Descripción} \\
  \midrule
  Diario    & Registros de un día específico. Por defecto muestra el día de hoy. \\
  Semanal   & Registros de la semana actual (lunes a hoy). \\
  Mensual   & Registros del mes actual. \\
  Semestral & Registros desde el inicio del semestre activo hasta hoy. \\
  \bottomrule
\end{tabular}
\end{center}

Para ver períodos anteriores, usa las flechas de navegación que
aparecen junto al selector de período (por ejemplo, para ver la semana
pasada o el mes anterior).

\subsection{Tarjetas de resumen del período}

Al seleccionar un período, aparecen tarjetas con:
\begin{itemize}
  \item Total de ingresos registrados en el período.
  \item Cantidad de personas únicas que ingresaron.
  \item Pico de ingresos (hora o día con más actividad).
\end{itemize}

\subsection{Gráfica de ingresos}

Debajo de las tarjetas hay una gráfica de barras que muestra la
distribución temporal de los ingresos:
\begin{itemize}
  \item \textbf{Diario}: Barras por hora del día.
  \item \textbf{Semanal}: Una barra por día de la semana.
  \item \textbf{Mensual}: Una barra por día del mes.
  \item \textbf{Semestral}: Una barra por mes.
\end{itemize}

\subsection{Tabla de registros}

Bajo la gráfica aparece la tabla detallada con todos los registros
del período seleccionado. Cada fila muestra:
\begin{itemize}
  \item Nombre e ID del usuario.
  \item Rol (estudiante, empleado, contratista).
  \item Fecha y hora exacta del ingreso.
\end{itemize}

\subsection{Exportar a Excel}

Para exportar los datos del período actual a un archivo Excel:

\begin{enumerate}
  \item Asegúrate de tener seleccionado el período que deseas exportar.
  \item Presiona el botón \textbf{Exportar Excel} (ícono de descarga).
  \item El navegador descargará automáticamente un archivo
        \texttt{.xlsx} con todos los registros visibles en pantalla.
\end{enumerate}

\begin{tip}
  El archivo Excel exportado se puede abrir directamente en
  Microsoft Excel, Google Sheets o LibreOffice Calc. Las columnas
  ya vienen con los encabezados en español.
\end{tip}

% ===============================================================
\section{Sección: Administradores}
\label{sec:admins}
% ===============================================================

\begin{nota}
  Esta sección solo es visible para usuarios con nivel
  \textbf{Superadmin}. Los administradores regulares no pueden
  acceder a ella.
\end{nota}

Esta sección permite gestionar las cuentas de acceso al panel de
administración. Desde aquí puedes crear nuevos administradores y
eliminar los existentes.

\subsection{Listado de administradores}

Muestra una tabla con todos los administradores registrados:
\begin{itemize}
  \item Nombre completo.
  \item ID institucional.
  \item Nivel (\textbf{Superadmin} en color dorado, \textbf{Admin} en azul).
  \item Fecha de creación.
  \item Fecha del último ingreso al panel.
\end{itemize}

\subsection{Crear un nuevo administrador}

\begin{enumerate}
  \item Haz clic en el botón \textbf{Nuevo administrador}.
  \item Se abre un formulario. Completa los campos:
    \begin{itemize}
      \item \textbf{ID institucional}: El número de identificación que
            usará para iniciar sesión.
      \item \textbf{Nombre completo}: Nombre y apellidos.
      \item \textbf{Nivel}: Elige \emph{Admin} o \emph{Superadmin}.
      \item \textbf{Contraseña}: Mínimo 8 caracteres.
      \item \textbf{Confirmar contraseña}: Debe coincidir con la anterior.
    \end{itemize}
  \item Presiona \textbf{Crear administrador}.
\end{enumerate}

\begin{aviso}
  Asigna el nivel \textbf{Superadmin} solo a personas de total confianza,
  ya que tendrán la capacidad de crear y eliminar otros administradores,
  incluyendo tu propia cuenta.
\end{aviso}

\subsection{Eliminar un administrador}

\begin{enumerate}
  \item Ubica al administrador en el listado.
  \item Presiona el ícono de eliminar (papelera) en su fila.
  \item Confirma la acción en la ventana emergente.
\end{enumerate}

\begin{peligro}
  No es posible eliminar la propia cuenta con la que estás iniciado
  sesión. Tampoco se puede dejar el sistema sin ningún superadmin.
\end{peligro}

\subsection{Cambiar la contraseña}

Actualmente el cambio de contraseña de un administrador se realiza
eliminando la cuenta y creándola nuevamente con la nueva contraseña.
Anota siempre las credenciales en un lugar seguro.

% ===============================================================
\section{Cerrar sesión}
% ===============================================================

Para cerrar sesión de forma segura:

\begin{enumerate}
  \item En la barra lateral, desplázate hasta el final.
  \item Presiona el botón \textbf{Cerrar sesión} (ícono de salida).
  \item El sistema te redirige automáticamente a la pantalla de login.
\end{enumerate}

\begin{aviso}
  Si cierras la pestaña del navegador sin cerrar sesión, tu sesión
  permanece activa. Cualquier persona que abra el panel desde ese mismo
  navegador podría acceder sin contraseña. Siempre cierra sesión cuando
  termines de trabajar, especialmente en equipos compartidos.
\end{aviso}

% ===============================================================
\section{Preguntas frecuentes}
% ===============================================================

\subsection*{¿El sistema funciona desde el celular?}

Sí. Tanto el portal de usuario como el panel de administración están
optimizados para funcionar en cualquier dispositivo (celular, tableta
o computador) siempre que esté conectado a la red WiFi del campus.

\subsection*{¿Qué pasa si el sistema no carga en el celular?}

Verifica que el celular esté conectado al WiFi del campus (no a datos
móviles). Si estás conectado al WiFi y no carga, el sistema puede
estar apagado: avisa al personal responsable de encenderlo.

\subsection*{¿Puedo usar el sistema desde fuera del campus?}

No. El sistema solo funciona dentro de la red WiFi institucional del
campus. No es accesible desde internet.

\subsection*{¿Por qué mi usuario aparece como bloqueado?}

El bloqueo puede deberse a una situación disciplinaria (reporte TIP)
o a una restricción administrativa. Debes contactar directamente a la
oficina encargada para conocer el motivo y solicitar la habilitación.

\subsection*{¿Puedo registrar mi ingreso varias veces en el mismo día?}

Sí. El sistema registra cada escaneo como un ingreso independiente.
Esto permite llevar un historial completo de entradas.

\subsection*{¿Los datos del semestre anterior se borran al cargar el nuevo?}

Los registros históricos de ingresos \textbf{no se borran}. Lo que
cambia es la lista de usuarios activos y su información académica o
laboral. El historial de accesos queda disponible en la sección de
Informes.

\subsection*{¿Qué es un reporte TIP?}

TIP (Tratamiento Institucional de Problemáticas) es el sistema
disciplinario de la universidad. Un reporte TIP registra una situación
que puede afectar el acceso del usuario. El Reglamento TIP se puede
descargar desde el portal de usuario.

% ===============================================================
\section{Glosario}
% ===============================================================

\begin{description}[leftmargin=3.5cm, labelwidth=3cm, style=nextline]
  \item[ID institucional] Número único que la universidad asigna a cada
    estudiante, empleado o contratista al momento de su vinculación.
  \item[Portal de usuario] Interfaz web accesible mediante código QR,
    usada para registrar ingresos.
  \item[Panel de administración] Interfaz web exclusiva para
    administradores del sistema.
  \item[Semestre activo] Período académico vigente configurado en el sistema.
  \item[Rol] Categoría del usuario: estudiante, empleado o contratista.
  \item[Estado] Condición de acceso del usuario: activo (puede ingresar)
    o bloqueado (no puede ingresar).
  \item[Reporte TIP] Registro disciplinario vinculado a un usuario.
  \item[CSV] Formato de archivo de texto separado por comas, usado para
    cargar listas masivas de usuarios. Se puede crear y editar con Excel.
  \item[Excel (.xlsx)] Formato de archivo de hoja de cálculo de Microsoft,
    generado al exportar los informes.
  \item[Superadmin] Nivel máximo de administrador, con permisos para
    gestionar otros administradores.
  \item[Admin] Nivel estándar de administrador, sin acceso a la gestión
    de otros administradores.
\end{description}

% ===============================================================
\vspace{2cm}
\begin{center}
  {\color{uccazul}\rule{0.5\textwidth}{0.5pt}}\\[0.5em]
  {\small\color{uccgris}
    Sistema de Control de Acceso UCC — Manual de Uso v1.0 \\
    Corporación Universitaria Comfacauca, 2026
  }
\end{center}

\end{document}
